% ===================================================================
%  CUADRO N.º 15 — El Mercado. — Los Alimentos.
%  Traduction 2026 d'apres texto/io, controlee sur texto/fr et
%  texto/en. Consigne : voir texto/es/10-cuadro-01.tex.
% ===================================================================

%%K t15-apar-1 apar
\VUsaut{4.0mm}
\VUtitre{15.0pt}{60}{CUADRO N.º 15}
\VUsaut{3.0mm}

%%K t15-tit-1 sub
\VUpk{11.4pt}{40}{El Mercado. --- Los Alimentos.}
\VUsaut{3.0mm}

%%K t15-01-1 p
1. --- La mayor parte de los comerciantes que nos proveen de los alimentos
necesarios están reunidos bajo la \VUgras{nave} \textsuperscript{(1)} del
\VUgras{mercado cubierto}, o en las calles vecinas.

%%K t15-02-1 p
2. --- El \VUgras{chacinero} \textsuperscript{(2)}, que vende \VUgras{jamón}
\textsuperscript{(3)}, \VUgras{morcilla} \textsuperscript{(4)}, \VUgras{salchichón}
\textsuperscript{(5)}, \VUgras{longaniza de tripa} \textsuperscript{(6)} y \VUgras{tocino}
\textsuperscript{(7)}, está al lado de la \VUgras{mantequera y quesera}
\textsuperscript{(8)}, cuyo puesto está cubierto de \VUgras{quesos de bola}
\textsuperscript{(9)} de corteza roja, \VUgras{camemberts} \textsuperscript{(10)}, grandes
\VUgras{ruedas de gruyer} \textsuperscript{(11)} y \VUgras{panes de mantequilla}
\textsuperscript{(12)} fresca.

%%K t15-03-1 p
3. --- Delante de ella, una \VUgras{verdulera} \textsuperscript{(13)} ofrece a sus
clientes hermosos \VUgras{tomates} \textsuperscript{(14)}, \VUgras{coliflores}
\textsuperscript{(15)}, \VUgras{lechugas} \textsuperscript{(16)}, \VUgras{coles}
\textsuperscript{(17)}, \VUgras{calabazas} \textsuperscript{(19)}, \VUgras{melones}
\textsuperscript{(18)} e incluso \VUgras{setas} \textsuperscript{(20)}. Pero un
\VUgras{cervecero}, que ha parado justo delante de su puesto su
\VUgras{carro de reparto} \textsuperscript{(21)} cargado de \VUgras{barriles de
cerveza} \textsuperscript{(22)}, impide que sus clientes se acerquen. Un hortelano
le trae una cesta de \VUgras{alcachofas} \textsuperscript{(23)}.

%%K t15-tit-2 sub
\VUsaut{4.0mm}
\VUpk{10.2pt}{40}{Vender y comprar.}
\VUsaut{3.0mm}

%%K t15-04-1 p
4. --- En el mercado la animación es grande, pues ha llegado la hora de la
\VUgras{subasta} \textsuperscript{(24)}. Las \VUgras{amas de casa} \textsuperscript{(26)} que
vienen aquí a hacer la compra se detienen de paso delante del puesto de la
\VUgras{casquera} \textsuperscript{(25)} para comprar \VUgras{callos} o una
\VUgras{cabeza de ternera}.

%%K t15-05-1 p
5. --- Un \VUgras{cocinero} \textsuperscript{(27)} grueso regatea un hermosísimo
\VUgras{atún} con la \VUgras{pescadera} \textsuperscript{(28)}, que sube los precios a
su antojo. Ha recibido una gran partida de \VUgras{anguilas, lenguados,
truchas} y \VUgras{carpas}. Su \VUgras{bacalao} y sus \VUgras{mejillones}
están muy frescos.

%%K t15-tit-3 sub
\VUsaut{4.0mm}
\VUpk{10.2pt}{40}{Lo barato y lo caro.}
\VUsaut{3.0mm}

%%K t15-06-1 p
6. --- La \VUgras{pollera} \textsuperscript{(29)}, instalada a su lado, es muy honrada.
Cuando vende un \VUgras{pavo}, una \VUgras{oca}, un \VUgras{pato} o un
simple \VUgras{pollo}, no pide más que el precio justo.

%%K t15-07-1 p
7. --- En la esquina de la plaza, en el primer piso, una
\VUgras{lencera} \textsuperscript{(30)} se ha establecido hace poco, donde los
compradores están siempre seguros de encontrar un surtido muy bueno de
\VUgras{manteles}, \VUgras{servilletas, delantales} y \VUgras{paños de
cocina}. En el mismo piso, un \VUgras{cuchillero} \textsuperscript{(31)} ha montado su
taller. Afila los \VUgras{cuchillos} de las verduleras y hace a los
hortelanos \VUgras{podaderas} del \VUgras{acero} más duro.

%%K t15-tit-4 sub
\VUsaut{4.0mm}
\VUpk{10.2pt}{40}{La tienda de comestibles.}
\VUsaut{3.0mm}

%%K t15-08-1 p
8. --- Debajo de su taller se ha abierto no hace mucho una gran tienda de
comestibles. Ya no es la \VUgras{tienda de ultramarinos} \textsuperscript{(32)} sencilla
y modesta de otros tiempos, que sólo vendía géneros corrientes ---
\VUgras{té} \textsuperscript{(33)}, \VUgras{chocolate} \textsuperscript{(34)},
\VUgras{pan de azúcar} \textsuperscript{(37)} y \VUgras{jabón} \textsuperscript{(38)}. Aquí se
vende de todo: \VUgras{galletas crujientes}, \VUgras{conservas}
\textsuperscript{(35)}, \VUgras{licores} \textsuperscript{(36)}, \VUgras{ron, coñac, kirsch,
anís} y \VUgras{curasao}. Se encuentra aquí caza con la misma facilidad ---
\VUgras{liebre} \textsuperscript{(39)}, \VUgras{tordos} \textsuperscript{(40)},
\VUgras{perdices} \textsuperscript{(41)}, \VUgras{codornices} \textsuperscript{(42)},
\VUgras{patos silvestres} \textsuperscript{(43)} o \VUgras{faisán} \textsuperscript{(44)} ---
que fruta tropical como \VUgras{plátanos} \textsuperscript{(46)} y \VUgras{piñas}.

%%K t15-09-1 p
9. --- Un \VUgras{dependiente} \textsuperscript{(45)} muy servicial está dispuesto a
entregar lo que se le pida: \VUgras{mermelada} \textsuperscript{(47)}, de grosella o
de membrillo, \VUgras{manteca} \textsuperscript{(48)}, \VUgras{pepinillos}
\textsuperscript{(49)} o \VUgras{sardinas} \textsuperscript{(50)} en salazón.

%%K t15-09-2 p
Eso es muy cómodo para las \VUgras{clientas} \textsuperscript{(51)}, que encuentran,
junto a los géneros más corrientes, las primicias más raras y la fruta más
sabrosa: \VUgras{peras} \textsuperscript{(52)} jugosas, \VUgras{ciruelas}
\textsuperscript{(53)}, \VUgras{uvas} \textsuperscript{(54)}, \VUgras{melocotones}
\textsuperscript{(55)}, \VUgras{almendras} \textsuperscript{(56)} y \VUgras{nueces}
\textsuperscript{(57)}.

%%K t15-tit-5 sub
\VUsaut{4.0mm}
\VUpk{10.2pt}{40}{Los clientes.}
\VUsaut{3.0mm}

%%K t15-10-1 p
10. --- El \VUgras{arroz} \textsuperscript{(58)} y las \VUgras{lentejas}
\textsuperscript{(59)} están junto a la caja de \VUgras{arenques} \textsuperscript{(60)}
ahumados; la \VUgras{patata} \textsuperscript{(61)} y la \VUgras{judía}
\textsuperscript{(63)} se codean con las \VUgras{manzanas} \textsuperscript{(62)}, los
\VUgras{higos} \textsuperscript{(64)}, las \VUgras{ciruelas pasas} \textsuperscript{(65)} y
las \VUgras{avellanas} \textsuperscript{(66)}. También se encuentran en esta tienda
\VUgras{huevos} \textsuperscript{(67)} frescos y \VUgras{aceitunas} \textsuperscript{(68)},
de modo que las \VUgras{cestas} \textsuperscript{(70)} y las \VUgras{redes de la
compra} \textsuperscript{(73)} se llenan a toda prisa.

%%K t15-11-1 p
11. --- El \VUgras{café} \textsuperscript{(72)} de esta excelente casa ha ganado una
fama bien merecida entre las \VUgras{criadas} \textsuperscript{(69)} y las cocineras,
pues el viejo \VUgras{tendero} \textsuperscript{(71)} lo tuesta él mismo con el mayor
esmero.

%%K t15-tit-6 sub
\VUsaut{4.0mm}
\VUpk{10.2pt}{40}{Cosas que ver.}
\VUsaut{3.0mm}

%%K t15-12-1 p
12. --- No todo el mundo viene al mercado a comprar provisiones. En el
pintoresco ir y venir de la callejuela que lleva a la nave, se ve a la
gente más variada. Junto a un bonachón \VUgras{matarife} \textsuperscript{(75)} que
empuja delante de sí una \VUgras{carretilla} \textsuperscript{(74)}, hay dos
soldados de permiso, muy orgullosos de lucir sus vistosos uniformes: el
\VUgras{húsar} \textsuperscript{(76)} con su \VUgras{dormán} azul y su
\VUgras{chacó de penacho}, y el \VUgras{cabo de zuavos} con sus calzones
anchos y un \VUgras{fez} en la cabeza.

%%K t15-13-1 p
13. --- El \VUgras{panadero} \textsuperscript{(80)}, de pie en el umbral de la
\VUgras{panadería} \textsuperscript{(78)}, acaba de amasar la masa de su
\VUgras{pan} \textsuperscript{(79)} y de sacar del horno los \VUgras{cruasanes}
\textsuperscript{(81)}, los \VUgras{panecillos} \textsuperscript{(82)} y las \VUgras{hogazas}
\textsuperscript{(83)} que están en el escaparate.

%%K t15-14-1 p
14. --- Encima de la panadería vive una hábil \VUgras{costurera}
\textsuperscript{(84)} que emplea a varias \VUgras{bordadoras} \textsuperscript{(85)}. A su
lado vive una \VUgras{lavandera y planchadora} \textsuperscript{(86)}, que almidona la
\VUgras{ropa blanca} \textsuperscript{(88)} después de lavarla y secarla, y la
plancha con una \VUgras{plancha} \textsuperscript{(87)}.

%%K t15-tit-7 sub
\VUsaut{4.0mm}
\VUpk{10.2pt}{40}{Carne, pescado y verduras.}
\VUsaut{3.0mm}

%%K t15-15-1 p
15. --- En la \VUgras{carnicería} \textsuperscript{(89)} de la planta baja se vende
carne de todas clases --- \VUgras{carnero} \textsuperscript{(90)}, \VUgras{vaca}
\textsuperscript{(94)} o \VUgras{ternera} \textsuperscript{(92)} --- en \VUgras{piernas de
cordero, filetes, asados} y \VUgras{chuletas}. El \VUgras{mancebo de
carnicería} \textsuperscript{(93)} corta toda esta carne y aparta los
\VUgras{huesos de tuétano} \textsuperscript{(96)}. A la puerta de la carnicería hay
una \VUgras{ostrera} \textsuperscript{(97)} que vende \VUgras{ostras} \textsuperscript{(98)}.
Las abre si se le pide.

%%K t15-16-1 p
16. --- Todos estos comerciantes pagan una patente o un derecho de puesto;
por eso el \VUgras{guardia} \textsuperscript{(100)} persigue sin piedad a las pobres
\VUgras{verduleras ambulantes} \textsuperscript{(99)}, que les hacen la competencia
llevando en sus carretillas \VUgras{zanahorias, rábanos, nabos, puerros} o
\VUgras{manojos de espárragos, guisantes frescos, espinacas, acederas,
berros} y \VUgras{perejil}.

%%K t15-tit-8 sub
\VUsaut{4.0mm}
\VUpk{10.2pt}{40}{¡Cuidado con el guardia!}
\VUsaut{3.0mm}

%%K t15-17-1 p
17. --- El muchacho que vende \VUgras{ajos} \textsuperscript{(101)} y \VUgras{cebollas}
sabe muy bien que a él tampoco le está permitido vender su mercancía cerca
del mercado. Echa a correr, a riesgo de volcar la cesta de
\VUgras{cangrejos}, \VUgras{cigalas} y \VUgras{quisquillas} de la
\VUgras{vendedora} \textsuperscript{(102)} de \VUgras{langostas} y
\VUgras{bogavantes}.

%%K t15-18-1 p
18. --- Todo el mundo se afana y arma un gran alboroto; pero eso no impide
oír claramente el reclamo del \VUgras{vidriero} \textsuperscript{(103)} ambulante,
ni el pregón del \VUgras{carbonero} \textsuperscript{(104)}, que acaba de dejar un
momento su carro para entregar un \VUgras{saco de carbón} \textsuperscript{(105)}.
\VUsaut{6.0mm}
\VUfilet{22mm}
